\chapter*{Introduction}
\addcontentsline{toc}{chapter}{Introduction}

When an equity option is priced with the Black--Scholes formula, a single parameter is not directly observable: the volatility $\sigma$. The \emph{implied volatility} of a quoted option is the value of $\sigma$ that reproduces its market price, and when computed across strikes $K$ and expiries $\tau$ it traces out the \emph{implied volatility surface} $(k,\tau)\mapsto \iv(k,\tau)$, with $k=\log(K/F)$ the log-moneyness against the forward. The surface is the market's summary of the risk-neutral distribution of the underlying at every horizon; it is the input to exotic pricing, hedging and risk systems, and its shape---the skew, the term structure---is economically informative in its own right.

The market, however, never exhibits a surface. It exhibits a \emph{sparse, noisy, irregular} cloud of quotes: only certain strikes and expiries trade, liquidity concentrates near the money, and bid--ask spreads widen rapidly in the wings. Reconstructing a complete, smooth surface from these observations is an ill-posed inverse problem: infinitely many surfaces interpolate the same quotes. The discipline that selects an admissible one is \emph{absence of static arbitrage}: the reconstructed prices must not allow riskless profit through butterfly or calendar strategies. These constraints admit a precise analytical form---non-negativity of the Durrleman function $g$ for butterfly arbitrage, monotonicity of total variance in maturity for calendar arbitrage---which makes them both testable and, crucially for this thesis, \emph{trainable}.

Three generations of methods address the problem. Parametric families, culminating in SVI and SSVI \cite{gatheral2014arbitrage}, offer interpretability and provable no-arbitrage guarantees at the cost of rigidity. Per-surface neural smoothers \cite{ackerer2020deep, chataigner2020deep, hoshisashi2024noarbitrage} recover flexibility by penalising arbitrage through the network's automatically-differentiated derivatives, but require one model per surface. Neural operators \cite{gonon2024operator} learn the smoothing \emph{map} itself: trained once across many days, a single discretisation-invariant model smooths any quote set without retraining.

\paragraph{Contributions.} This thesis makes four contributions.
\begin{enumerate}
  \item \textbf{A unified benchmark.} All three generations are implemented and evaluated on the \emph{same} daily S\&P~500 dataset (OptionMetrics, 2018--2022), under the \emph{same} protocol: hold-out reconstruction error, a dense-grid static-arbitrage audit, robustness to quote sparsity, and per-regime analysis. Published papers each use their own data and metrics; to our knowledge no like-for-like comparison exists.
  \item \textbf{A validated replication.} The complete Gatheral--Jacquier toolkit---three SVI parameterisations and their exact conversions, the butterfly-elimination algorithm, the calibration recipes, the SSVI no-arbitrage theorems, arbitrage-free interpolation---is replicated and validated against the paper's own numerical examples, correcting in passing a typographical omission in the printed statement of the Jump-Wings inversion (Lemma~3.2).
  \item \textbf{A systematic arbitrage audit.} Beyond fit error, every method is scored on \emph{how much} static arbitrage it leaves on real data---violation rates and magnitudes of the Durrleman condition on dense grids, cross-slice crossedness, and term-structure inversions---quantities the literature rarely reports.
  \item \textbf{The low-data operator experiment.} Operator deep smoothing was demonstrated on $\sim$60 million intraday quotes; academic data access is daily, two orders of magnitude smaller. We measure operator performance in this regime and test whether pre-training on synthetic SSVI surfaces closes the gap (regimes R1/R2/R3: real-only, synthetic-only, pre-train-then-fine-tune).
\end{enumerate}

\paragraph{Outline.} Chapter~\ref{ch:foundations} sets the financial and mathematical foundations. Chapter~\ref{ch:literature} reviews the three methodological generations. Chapter~\ref{ch:data} documents the data pipeline. Chapters~\ref{ch:parametric}--\ref{ch:operator} develop the three method families, each validated on synthetic ground truth before being applied to market data. Chapter~\ref{ch:evaluation} performs the unified comparison, and the conclusion discusses limitations and extensions.
